\documentclass[12pt, twoside]{article}
\input{/Users/chris/github/course-files/docs/ib/preamble}

\fancyhead[LE]{\thepage}
\fancyhead[RO]{Name: \hspace{4cm} \,\\ \hspace{2.9cm} \,}
\fancyhead[LO]{La Scuola d'Italia / Dr.\ Huson / IB Math HL \\* 21 May 2026}

\begin{document}
\subsubsection*{5.1 HW: Maclaurin Series --- Take-home (calculator allowed)}

Two problems. Answer on lined paper. Show all working.
Total: 36 marks. A GDC may be used; show algebraic working before any
numerical evaluation.

\begin{enumerate}[itemsep=0.15cm]

%--- Problem 4: e^x sin x series, integral approx, ODE-style identity, limit [21 marks] ---
%--- Source: AA HL 2022 May P1 TZ1 Q12 ---
\item[4.] \makebox[\linewidth][s]{[Maximum mark: 21]\hfill {\small AA HL 2022 May P1 TZ1 Q12}}

The function $f$ is defined by $f(x) = e^{x} \sin x$, where
$x \in \mathbb{R}$.

\begin{enumerate}[itemsep=0.15cm]
  \item Find the Maclaurin series for $f(x)$ up to and including the
  $x^{3}$ term. \hfill [4]
  \item Hence, find an approximate value for
  $\displaystyle\int_{0}^{1} e^{x^{2}} \sin\!\left(x^{2}\right)\, dx$.
  \hfill [4]
\end{enumerate}

The function $g$ is defined by $g(x) = e^{x} \cos x$, where
$x \in \mathbb{R}$.

\begin{enumerate}[itemsep=0.15cm]
  \item[(c)] \begin{enumerate}[itemsep=0.15cm]
    \item[(i)] Show that $g(x)$ satisfies the equation
    $g''(x) = 2\!\left(g'(x) - g(x)\right)$. \hfill [4]
    \item[(ii)] Hence, deduce that
    $g^{(4)}(x) = 2\!\left(g'''(x) - g''(x)\right)$. \hfill [1]
  \end{enumerate}
  \item[(d)] Using the result from part (c), find the Maclaurin series
  for $g(x)$ up to and including the $x^{4}$ term. \hfill [5]
  \item[(e)] Hence, or otherwise, determine the value of
  \[
  \lim_{x \to 0}\frac{e^{x} \cos x - 1 - x}{x^{3}}.
  \]
  \hfill [3]
\end{enumerate}

\newpage

%--- Problem 5: Arctan series from geometric series; alternating-series error [15 marks] ---
%--- Source: AA HL 2025 Nov P3 TZ3 Q02 (parts a-e only; calc parts omitted) ---
\item[5.] \makebox[\linewidth][s]{[Maximum mark: 15]\hfill {\small AA HL 2025 Nov P3 TZ3 Q02 (a)--(e)}}

The following problem uses Maclaurin series to investigate approximations
of mathematical constants and the accuracies of such approximations.

\begin{enumerate}[itemsep=0.15cm]
  \item Given $|x| < 1$, find the sum to infinity of the geometric
  series
  \[
  1 - x^{2} + x^{4} - x^{6} + \cdots.
  \]
  \hfill [2]
  \item Hence, use integration to show that the Maclaurin series of
  $\arctan x$ may be expressed as
  \[
  \arctan x = x - \frac{x^{3}}{3} + \frac{x^{5}}{5} - \frac{x^{7}}{7}
  + \cdots.
  \]
  \hfill [3]
  \item Using $x = \dfrac{1}{\sqrt{3}}$ and the first \textbf{three}
  (non-zero) terms of the Maclaurin series of $\arctan x$, find an
  approximation for $\pi$ to three decimal places. \hfill [3]
\end{enumerate}

The Maclaurin series of $\arctan x$ is an alternating series.

\textbf{Theorem.} For an alternating series whose terms decrease in
magnitude, the error obtained in using a finite number of terms is less
than or equal to the absolute value of the next term in the sequence.

By this theorem, the maximum error in using the first three (non-zero)
terms as an approximation to $\arctan x$ is $\left|-\dfrac{x^{7}}{7}\right|$,
that is,
\[
\left|\arctan x - \left(x - \frac{x^{3}}{3} + \frac{x^{5}}{5}\right)\right|
\leq \left|-\frac{x^{7}}{7}\right|.
\]

\begin{enumerate}[itemsep=0.15cm]
  \item[(d)] Determine how many (non-zero) terms of the series would
  need to be used so that the error in approximating
  $\arctan\!\left(\dfrac{1}{\sqrt{3}}\right)$ is less than $0.0001$.
  \hfill [3]
  \item[(e)] By using integration by parts, show that
  \[
  \int_{0}^{1/\sqrt{3}} \arctan x\, dx
  = \frac{\pi}{6\sqrt{3}} - \frac{1}{2}\ln\frac{4}{3}.
  \]
  \hfill [4]
\end{enumerate}

\end{enumerate}
\end{document}
